% Kamp III.2, Lemma 10 and the finite Boolean normalization used in its final application.

\subsubsection{All regional interval components}
\label{sec:p09-kamp-iii2-lemma10}

The recursive \(\operatorname{Bet}\) classes use an existential
conjunction or a universal disjunction over a strict interior point.
Lemmas~5 and 9 handle those two constructors. Lemma~10 therefore
provides the interval property for every finite stage, and conjunction
adds the endpoint literals needed in the concluding proof
\cite[ch.~III, lemma~10, p.~93]{Kamp1968}.

\begin{hbmethod}{def:p09-kamp-bet-construction}{The stage construction for regional interval formulas}
Use the exact unstarred recursion of
\cref{def:p09-kamp-iii1-components}. Stage zero is empty.
At a retained component reuse its earlier representation. At an
existential constructor, separate the finite matrix into its At
conditions at the quantified point and its two Bet families on either
side. Apply Lemma~5 to their supplied representations. At a universal
constructor apply Lemma~9 to those three disjunction families.

The At conditions have elementary histories; the constituent Bet
conditions come from earlier stages. Keep the strict guard
\(x<z<y\), including when a matrix family is empty. For a
conjunction of finished Bet components and At conditions at \(x\)
or \(y\), insert each endpoint formula using the interval endpoint
lemma and apply the complete finite fusion construction to all factors.
An empty outer conjunction is represented by the one-interval truth word.
\end{hbmethod}

\begin{hblemma}{thm:p09-kamp-iii2-lemma10}{proved=annex:p9.kamp.iii2.lemma10}{Kamp's interval-component lemma}
Every member of \(\operatorname{Bet}(x,y)\) has property
\(I(x,y)\). Every finite conjunction of members of
\[
       \operatorname{Bet}(x,y)\cup\operatorname{At}(x)
                                      \cup\operatorname{At}(y)
\]
also has that property. For a supplied finite construction history,
the stated algorithm returns a finite disjunction of actual
\(\operatorname{Ep}(x,y)\) formulas, equivalent under \(x<y\)
over \(\mathsf{Com}\). This includes the declared empty-family
conventions and retains both source quantifier constructors
\cite[ch.~III, lemma~10, p.~93]{Kamp1968}.
\end{hblemma}

\begin{hbdefinition}{def:p09-kamp-interval-boolean}{Positive Boolean expressions in interval components}
An interval Boolean expression on \((x,y)\) is a finite Boolean
combination of members of
\(\operatorname{Bet}(x,y)\cup\operatorname{At}(x)\cup
\operatorname{At}(y)\). It is \emph{positive} when the external
Boolean constructors are only conjunction and disjunction. A negative
unary literal inside an At component remains permitted.
Use \(T^K_x\) and \(F^K_x\) for empty Boolean truth and
falsehood. No claim is made that these Boolean expressions themselves
are single members of \(\operatorname{Bet}\).
\end{hbdefinition}

\begin{hblemma}{thm:p09-kamp-interval-boolean-normalization}{proved=annex:p9.kamp.interval.boolean.normalization}{Finite interval negation normalization}
Every interval Boolean expression has an effectively constructed positive
interval Boolean expression equivalent to it over \(\mathsf{Lin}\).
In particular the negation of a Bet component can be replaced by such
a finite expression without adding a Bef or Aft component or an extra
free variable.

Consequently every interval Boolean expression has property \(I(x,y)\):
apply Lemma~10 and finite fusion to the positive expression. The
normalization supplies explicitly the negation and Boolean distribution
step used by Kamp before his final existential reduction
\cite[ch.~III, pp.~61--62, 93]{Kamp1968}.
\end{hblemma}

\begin{hbexample}{def:p09-kamp-bet-empty-interior}{Existence and universal vacuity are retained}
The stage-one formulas
\[
 \exists z\,(x<z<y\land T^K_z),\qquad
 \forall z\,(x<z<y\to F^K_z)
\]
assert respectively that the open interval has a point and that it
has no point. They are complements on every linear order. The first
has a two-interval truth word as an elementary representative; the
second has the one-interval word \((T^K,F^K,T^K)\).
At adjacent endpoints only the second holds. The normalization keeps
these strict witness and vacuity conditions when moving a negation
across a regional quantifier.
\end{hbexample}
